\documentclass{article}

\usepackage{lewis}
\usepackage{bohr}

\title{Chem 001A\\ Week 7}
\author{Danny Topete}
\date{November 10, 2025}

\begin{document}
\maketitle

\section{Monday Lecture}

\begin{itemize}
  \item Discussion this week will just go over the midterm
  \item No mastery exams in November, therefore testing center
  \item Please re-do unit 3 (gradescope)
  \item You need to sign up for unit 4 quiz, testing center fml
  \item Seriously, do it this week
  \item Midterm and finals are their own score, they can't be applied
    to mastery exams (figure out mastery/midterm works)
\end{itemize}

\subsection{Hybrid Orbitals in Valence Bond Theory}
\begin{itemize}
  \item Coming back to quantum and taking about hybrid orbitals
  \item Valence Bond Model
  \item Electrons are in an orbital
  \item Valence Bond model, describes covalent bond as overlap of half-filled atomic orbitals
  \item Attraction between negatively charged electron pair and two atom's pos charged nuclei physically links the two atoms forming a
    covalent bond
\end{itemize}

\subsection{Predicting Geometry with VB Model}
\begin{itemize}
  \item According to VB model, bonding $H_2S$ can be understood as an overlap between
    half-filled 1s orbital of two hydrogen atoms and two half-filled 3p orbitals on sulfur.

   H - 1s, H - 1s, S $3s^2$ + $3p^4$
\end{itemize}
\subsection{Atomic orbitals limit VB Theory}
\begin{itemize}
  \item H - 1s, C - $2s^2$ + $2p^2$, we get a really odd molecule where
    carbon doesn't have a full octet and hydrogen bonds have a 90* angle
\end{itemize}

\subsection{Meolecular geometry and hybrid orbitals}
\begin{itemize}
  \item Model BeCl$_2$, \lewis{Cl}{.}{.}{.}{.}{}{}{.}{.}  - Be - \lewis{Cl}{}{}{.}{.}{.}{.}{.}{.} (linear)
\end{itemize}

\begin{enumerate}
  \item Hybrid orbitals don't exist in isolated atoms. Formed only covalently
  \item shape of hybrid orbitals depend on type and num of atomic orbitals used in their construction
  \item Set of hybrid orbitals constructed by combining a set of atmic orbitals. N orbitals each set is always the same
  \item All orbitals in set hybrid orbitals equivalent in shape and energy but orient in different directions.
  \item Type hybrid orbitals formed in bonded atoms
\end{enumerate}

\pagebreak

\section{Friday Lecture 11/14}
\subsection{Molecular Mass and molar mass}
\begin{itemize}
  \item Practice problem
    How many atoms in 0.852 moles of

    Fe$^{+3}O_3^{-2}$ \times{} NA \times{} $\frac{5 atoms}{per molecule of Fe^{+3}O_3^{-2}} $
    = 2.5654\cdot{}$10^{24}$ atoms
\end{itemize}

\subsection{Percentage Composition}
\begin{itemize}
  \item Percentage composition
    $$ \%A = \frac{Mass{} A}{mass{} compound} \times{} 100\% $$
  \item Chemical formula tells you the \% composition
\end{itemize}

\end{document}
